% !PRINTED 60
\weberpage{60}

\begin{center}
\S.\,16.

\gesperrt{Die Gruppen der Zahlclassen nach einem zusammengesetzten Modul}.
\end{center}

Aus der Verbindung der beiden Sätze 2. und 3. der beiden vorangegangenen Paragraphen ergiebt sich nun folgendes Resultat, durch welches die Aufgabe, die wir am Anfang des \S.\,14 gestellt haben, vollständig gelöst wird:

\noindent\textbf{4.} \gesperrt{Wenn der Modul}

\[
m = 2^\lambda q_1^{\varkappa_1} q_2^{\varkappa_2} \ldots
\]

\gesperrt{ist, und} $g_1, g_2 \ldots$ \gesperrt{primitive Wurzeln der Quadrate der Primzahlen} $q_1, q_2 \ldots$ \gesperrt{sind, wenn ferner}

\[
a = 2, \quad b = \tfrac12 \varphi(2^\lambda), \quad c_1 = \varphi(q_1^{\varkappa_1}), \quad c_2 = \varphi(q_2^{\varkappa_2}) \ldots
\]

\gesperrt{ist, so kann man für jede Zahl} $n$, \gesperrt{die zu} $m$ \gesperrt{relativ prim ist, ein System von Zahlen} $\alpha, \beta, \gamma_1, \gamma_2 \ldots$ \gesperrt{nach den Moduln} $a, b, c_1, c_2 \ldots$ \gesperrt{eindeutig bestimmen, die den Gleichungen}

\weq{1}{
\begin{aligned}
n &\equiv (-1)^\alpha 5^\beta \pmod{2^\lambda} \\
&\equiv g_1^{\gamma_1} \pmod{q_1^{\varkappa_1}} \\
&\equiv g_2^{\gamma_2} \pmod{q_2^{\varkappa_2}} \\
&\; \cdots\cdots\cdots\cdots\cdots\cdots
\end{aligned}
}

\gesperrt{genügen.}

\gesperrt{Die Zahlen $\alpha, \beta, \gamma_1, \gamma_2 \ldots$ heissen das System der Indices von $n$ für den Modul $m$.}

Hier ist zunächst $\lambda \geqq 3$ vorausgesetzt. Der Satz gilt aber auch für die übrigen Werthe von $\lambda$, wenn

\[
\text{für } \lambda = 0, 1, \quad a = 1, \quad b = 1
\]
\[
\text{\textquotedbl{}} \quad \lambda = 2, \qquad\quad a = 2, \quad b = 1
\]

gesetzt wird. Für $\lambda = 0, 1$ sind die Indices $\alpha$ und $\beta = 0$ zu setzen oder auch ganz wegzulassen, für $\lambda = 2$ fällt $\beta$ weg, während $\alpha$ gleich 0 oder gleich 1 sein kann.

Um nun also die Gruppe $\mathfrak{N}$ der Zahlclassen $N$ nach dem Modul $m$ durch eine Basis darzustellen, bestimme man die Zahlclassen

\weq{2}{A, B, C_1, C_2 \ldots,}

oder wenigstens Repräsentanten dieser Classen, was immer möglich ist, aus den Congruenzen
%CONT>>
